% !TeX root = ../../Book5.tex
% 纲要标题：### A.4 Hecke 代数与 Kazhdan--Lusztig 理论
% 纲要位置：工作记录/计划纲要.md，第 4823 行。
% #### A.4.1 参数、二次关系与标准基
% #### A.4.2 反演算子与规范基
% #### A.4.3 胞腔与一个二维表示
\section{Hecke 代数与 Kazhdan--Lusztig 理论}
\label{b5:sec:A04}

群代数中简单反射满足平方为一。如果把这个二次关系变形，同时保留辫关系，就得到 Hecke 代数。它的标准基仍由群元素标号，却还存在与一个参数反演相容的特殊基。本节固定一种归一化，并通过 \(S_3\) 把不同基与胞腔表示具体算出来。

\subsection{参数、二次关系与标准基}
\label{b5:subsec:A04:01}

\begin{definition}\label{b5:def:hecke-algebra}
设 \((W,S)\) 为 Coxeter 系统，\(v\) 为形式变量，\(R=\mathbb Z[v,v^{-1}]\)。用生成元 \(H_s\)（\(s\in S\)）、关系
\[
(H_s-v)(H_s+v^{-1})=0
\]
以及与 Coxeter 系统相同的辫关系定义的含单位结合 \(R\) 代数，称为等参数\term[Hecke-daishu]{Hecke 代数}[Hecke algebra]，记为 \(\mathcal H(W)\)。
\end{definition}
\symidx{\(\mathcal H(W)\)：Coxeter 系统的等参数 Hecke 代数}

\begin{theorem}[Hecke 标准基]\label{b5:thm:hecke-standard-basis}
设 \((W,S)\) 为 Coxeter 系统，\(\ell\) 为关于 \(S\) 的长度函数，\(v\) 为形式变量，\(R=\mathbb Z[v,v^{-1}]\)。令 \(\mathcal H(W)\) 为 \(R\) 上的等参数 Hecke 代数，以 \(H_s\)（\(s\in S\)）为生成元，满足 \((H_s-v)(H_s+v^{-1})=0\) 及该 Coxeter 系统的辫关系。对约化词 \(w=s_1\cdots s_r\)，乘积
\(H_w=H_{s_1}\cdots H_{s_r}\) 与约化词的选择无关，而且 \(\{H_w\mid w\in W\}\) 为 \(R\) 基。乘法满足
\begin{equation}\label{b5:eq:hecke-standard-multiplication}
H_sH_w=
\begin{cases}
H_{sw},&\ell(sw)=\ell(w)+1,\\
H_{sw}+(v-v^{-1})H_w,&\ell(sw)=\ell(w)-1.
\end{cases}
\end{equation}
\end{theorem}
\symidx{\(H_w\)：Hecke 代数的标准基元}
\begin{proof}
记 \(d=v-v^{-1}\)。由\cref{b5:thm:coxeter-word-properties}的辫词性质，\(H_w\) 良定义。若 \(sw\) 较短，则 \(w=s(sw)\) 是约化分解，二次关系 \(H_s^2=1+dH_s\) 给出第二种乘法；较长时直接接上 \(s\)。因此所有生成元乘积都在这些 \(H_w\) 的生成空间中。

为证线性无关，取自由模 \(E=\bigoplus_{w\in W}Re_w\)，按同一长度规则定义算子
\[
 L_se_w=e_{sw}+\begin{cases}0,&\ell(sw)>\ell(w),\\de_w,&\ell(sw)<\ell(w),\end{cases}
 \quad
 R_te_w=e_{wt}+\begin{cases}0,&\ell(wt)>\ell(w),\\de_w,&\ell(wt)<\ell(w).\end{cases}
\]
在 \(e_w,e_{sw}\) 上直接算得 \(L_s^2=1+dL_s\)。为验证 \(L_sR_t=R_tL_s\)，记 \(n=\ell(w)\)，先设 \(sw\ne wt\)。当 \(\ell(sw),\ell(wt)\) 分别为 \(n+1,n-1\) 或 \(n-1,n+1\) 时，\(\ell(swt)\) 与这两个数都相差一，故只能为 \(n\)。当两者都为 \(n+1\) 时，若 \(\ell(swt)=n\)，对约化词 \(wt\) 左乘 \(s\) 应用交换性质。若删去 \(w\) 中的位置，再右乘 \(t\) 就得到长度至多为 \(n-1\) 的 \(sw\)，矛盾；若删去末尾的 \(t\)，则 \(swt=w\)，即 \(sw=wt\)，仍矛盾。因此此时 \(\ell(swt)=n+2\)。

当 \(\ell(sw)=\ell(wt)=n-1\) 时，若 \(\ell(swt)=n\)，令 \(z=sw\)。则 \(sz=w\)、\(zt=swt\) 都长 \(n\)，且二者不同。刚证的两次长度上升情形给出 \(\ell(szt)=n+1\)，而 \(szt=wt\) 仅长 \(n-1\)，矛盾。因此 \(\ell(swt)=n-2\)。将这四种长度关系代入算子定义，两种复合的共同值为
\[
\begin{array}{c|c|l}
(\ell(sw)-n,\ell(wt)-n)&\ell(swt)&L_sR_te_w=R_tL_se_w\\\hline
(1,1)&n+2&e_{swt}\\
(1,-1)&n&e_{swt}+de_{sw}\\
(-1,1)&n&e_{swt}+de_{wt}\\
(-1,-1)&n-2&e_{swt}+d(e_{sw}+e_{wt})+d^2e_w
\end{array}
\]
若 \(sw=wt\)，两次长度改变方向相同。都增加时，两种复合均为 \(e_w+de_{sw}\)；都减少时，均为 \((1+d^2)e_w+de_{sw}\)。于是左右算子确实交换。

当 \(m_{st}<\infty\) 时，\(s,t\) 的两条长度 \(m_{st}\) 交替词是二面体子群 \(W_{\{s,t\}}\) 最长元的约化词。由上一节所用的标准抛物子群长度相容性，它们在 \(W\) 中也约化，故相应左算子乘积在 \(e_1\) 上有相同值。两种乘积都与全部 \(R_u\) 交换，而各个 \(e_w\) 都能由 \(e_1\) 沿一个约化词依次右乘得到，故两种乘积在全 \(E\) 上相同。这样 \(L_s\) 满足全部定义关系，给出 \(\mathcal H(W)\) 在 \(E\) 上的表示。最后 \(H_we_1=e_w\)，所以任何 \(H_w\) 的线性关系作用到 \(e_1\) 后都只能有零系数。
\end{proof}

原始 Hecke 代数的基及乘法约定见\cite[Section 1, p. 165]{KazhdanLusztig1979Hecke}。原文使用 \(q=v^2\) 与 \(T_w=v^{\ell(w)}H_w\)，于是二次式变为 \((T_s+1)(T_s-q)=0\)。比较公式时须先写清参数与归一化的对应。

\ProseParagraphBreak
在秩一情形无需任何一般词理论：
\[
\mathcal H(C_2)=R[H]/\bigl(H^2-(v-v^{-1})H-1\bigr).
\]
首一除法直接给出基 \(1,H\)。当 \(v=1\) 时，关系成为 \(H^2=1\)，恢复群环 \(\mathbb Z[C_2]\)；一般 Coxeter 情形由标准基定理同样得到
\(\mathcal H(W)/(v-1)\simeq\mathbb Z[W]\)。
\ProseParagraphBreak
在 \(S_3=\langle s,t\rangle\) 中，标准基为
\[
1,\ H_s,\ H_t,\ H_{st},\ H_{ts},\ H_{sts},
\]
而 \eqref{b5:eq:hecke-standard-multiplication}给出
\[
H_sH_{st}=H_t+(v-v^{-1})H_{st},\qquad
H_sH_{ts}=H_{sts}.
\]
一个乘法用了长度下降，另一个用了长度上升；不能在两种情况下都直接把群元素下标相乘。

\subsection{反演算子与规范基}
\label{b5:subsec:A04:02}

二次关系给出 \(H_s^{-1}=H_s-(v-v^{-1})\)。由此可定义半线性环对合
\[
\overline v=v^{-1},\qquad
\overline{H_w}=H_{w^{-1}}^{-1}.
\]
它保持乘法而反演系数参数。良定义可从生成关系核验：逆元满足参数反演后的二次关系，取逆并反向读取辫关系仍得到同一对辫词；再作一次变换恢复各生成元。因此它不是把乘法次序倒转的反自同构。

\begin{theorem}[Kazhdan--Lusztig 规范基]\label{b5:thm:kl-basis}
设 \((W,S)\) 为 Coxeter 系统，\(\ell\) 为关于 \(S\) 的长度函数，\(v\) 为形式变量，\(R=\mathbb Z[v,v^{-1}]\)。令 \(\mathcal H(W)\) 为 \(R\) 上的等参数 Hecke 代数，满足 \((H_s-v)(H_s+v^{-1})=0\) 及 Coxeter 辫关系，\(H_w\) 为其标准基。以 \(\overline v=v^{-1}\)、\(\overline{H_w}=H_{w^{-1}}^{-1}\) 定义半线性环对合。对每个 \(w\in W\)，存在唯一元素 \(C'_w\in\mathcal H(W)\) 满足
\[
\overline{C'_w}=C'_w,\qquad
C'_w=H_w+\sum_{y<w}a_{y,w}(v)H_y,\qquad
a_{y,w}(v)\in v^{-1}\mathbb Z[v^{-1}],
\]
其中 \(<\) 为 Bruhat 严格序。这些元素构成 \(\mathcal H(W)\) 的 \(R\) 基，且对形式变量 \(q\) 有
\[
a_{y,w}(v)=v^{\ell(y)-\ell(w)}\cdot P_{y,w}(v^2),
\qquad P_{y,w}(q)\in\mathbb Z[q],
\]
对 \(y<w\) 有 \(\deg P_{y,w}\le\bigl(\ell(w)-\ell(y)-1\bigr)/2\)，并置 \(P_{w,w}=1\)。
\end{theorem}
\symidx{\(C'_w,P_{y,w}(q)\)：Kazhdan--Lusztig 规范基元与多项式}

\begin{proof}
先注意 \(\overline{H_w}=H_w+\sum_{y<w}r_{y,w}H_y\)：沿约化词把每个 \(H_s\) 换成 \(H_s-d\)，非最高项来自删去位置的子词，利用乘法规则和子词判据便得到所示三角形展开。由\cref{b5:thm:coxeter-word-properties}的第三项，主下区间 \([1,w]\) 有限，因此下面对固定 \(w\) 的各个和式均为有限和。

按长度构造 \(C'_w\)，取 \(C'_1=1\)。设 \(w=su\) 且 \(\ell(w)=\ell(u)+1\)，较短元素的规范基已经构造。乘积
\(D=(H_s+v^{-1})C'_u\) 自对偶。最高项 \(H_u\) 的贡献为 \(H_w+v^{-1}H_u\)；对 \(y<u\)，乘法规则给出
\[
(H_s+v^{-1})a_{y,u}H_y=
\begin{cases}
a_{y,u}H_{sy}+v^{-1}a_{y,u}H_y,&sy>y,\\
a_{y,u}H_{sy}+v a_{y,u}H_y,&sy<y.
\end{cases}
\]
子词判据保证这里的下标 \(y,sy\) 都小于 \(w\)。各个 \(a_{y,u}\) 只含严格负次幂，故上式的系数都没有正次幂；常数项只可能来自 \(sy<y\) 时的 \(v a_{y,u}\)，并等于 \(a_{y,u}\) 中 \(v^{-1}\) 的系数 \(\mu(y,u)\)。因此定义
\[
 C'_w=D-\sum_{\substack{y<u\\sy<y}}\mu(y,u)C'_y.
\]
减项仍自对偶，恰消去各个常数项；其低项系数原本只有负次幂，故不会产生新的常数项。这给出所需存在性。

同时归纳验证：\(a_{y,w}\) 中每个指数不小于 \(\ell(y)-\ell(w)\)，且与该数同奇偶。设 \(a_{y,u}\) 的一个指数为 \(k\)，则归纳假设给出
\[
k\ge\ell(y)-\ell(u),\qquad
k\equiv\ell(y)-\ell(u)\pmod2.
\]
在 \(sy>y\) 的情形中，\(H_{sy}\) 的指数 \(k\) 及 \(H_y\) 的指数 \(k-1\) 分别不小于 \(\ell(sy)-\ell(w)\) 及 \(\ell(y)-\ell(w)\)，奇偶也分别一致。在 \(sy<y\) 的情形中，\(H_{sy}\) 的指数 \(k\) 及 \(H_y\) 的指数 \(k+1\) 的下界分别比上述目标下界多二，故同样满足要求。最高项的额外贡献 \(v^{-1}H_u\) 也满足这一约束。

再检查减项。为统一记号，置 \(a_{y,y}=1\)。若 \(\mu(y,u)\ne0\)，则 \(-1\equiv\ell(y)-\ell(u)\pmod2\)，所以 \(\ell(w)-\ell(y)\) 为偶数。对每个 \(x\le y\)，\(C'_y\) 中 \(H_x\) 的系数的指数下界为
\[
\ell(x)-\ell(y)\ge\ell(x)-\ell(w),
\qquad
\ell(x)-\ell(y)\equiv\ell(x)-\ell(w)\pmod2.
\]
这也包括 \(x=y\) 时的指数零。因此减项保持所需的下界与奇偶性。结合 \(C'_w\) 的所有下位系数只含严格负次幂，便可唯一写成
\(v^{\ell(y)-\ell(w)}P_{y,w}(v^2)\)。若 \(P_{y,w}\) 中出现 \(q^j\)，则 \(\ell(y)-\ell(w)+2j\le-1\)，这正给出所述次数上界。

若有两个符合刻画的元素，取其差中 Bruhat 极大的非零下标 \(y\)。三角形反演说明对应系数 \(f\) 必满足 \(f(v)=f(v^{-1})\)。但 \(f\) 只有严格负次幂，只能为零，矛盾。这证明唯一性。最后，对每个 \(w\)，在有限主下区间 \([1,w]\) 上反解单位三角形变换，即可把 \(H_w\) 写成有限个 \(C'_y\) 的线性组合；同一三角形关系也给出线性无关性，故全部 \(C'_w\) 组成 \(R\) 基。
\end{proof}

这是原论文\cite[Theorem 1.1, equivalent formulation (1.1.c), p. 166]{KazhdanLusztig1979Hecke}的归一化改写。多项式 \(P_{y,w}\) 称为\term[Kazhdan-Lusztig-duoxiangshi]{Kazhdan--Lusztig 多项式}[Kazhdan--Lusztig polynomial]。下述低秩计算可直接检验递推；一般系数的正性以及 Lie 代数中的特征标公式则是另外的结果。

\ProseParagraphBreak
对简单反射，
\[
C'_s=H_s+v^{-1}.
\]
由逆元公式，\(\overline{H_s+v^{-1}}=H_s-(v-v^{-1})+v=H_s+v^{-1}\)，所以确为自对偶。若先只要求 \(H_s+a(v)\) 自对偶，条件其实是
\[
a(v)-a(v^{-1})=v^{-1}-v.
\]
除了解 \(a(v)=v^{-1}\)，还可加上任意满足 \(b(v)=b(v^{-1})\) 的 Laurent 多项式。要求 \(a(v)\) 只含严格负次幂，便排除了这份自由度：两解的差同时在参数反演下不变且只含负次幂，只能为零。这正是次数条件的规范化作用。若记 \(a=v^{-1}\)，在 \(S_3\) 中得到
\begin{align*}
C'_{st}&=C'_sC'_t
=H_{st}+a(H_s+H_t)+a^2,\\
C'_{ts}&=C'_tC'_s
=H_{ts}+a(H_s+H_t)+a^2,\\
C'_{sts}&=C'_sC'_tC'_s-C'_s\\
&=H_{sts}+a(H_{st}+H_{ts})+a^2(H_s+H_t)+a^3.
\end{align*}
每个乘积和差均自对偶；展开后最高项系数一，其余系数只含严格负次幂，且下标在已计算的 Bruhat 区间中。因此它们就是所求规范基。最后一个式子的减项不可省略：未减之前 \(H_s\) 的系数为 \(1+a^2\)，含常数项，不满足规范条件。这里所有 \(P_{y,w}\)（\(y\le w\)）都等于一，但高秩时不会一直如此。

\subsection{胞腔与一个二维表示}
\label{b5:subsec:A04:03}

\begin{definition}\label{b5:def:kl-left-cell}
设 \((W,S)\) 为 Coxeter 系统，\(v\) 为形式变量，\(\mathcal H(W)\) 为 \(\mathbb Z[v,v^{-1}]\) 上的等参数 Hecke 代数，满足 \((H_s-v)(H_s+v^{-1})=0\)，\(C'_w\)、\(w\in W\)，为 Kazhdan--Lusztig 规范基。对 \(x,y\in W\)，若某个 \(s\in S\) 使 \(C'_sC'_y\) 的展开含非零的 \(C'_x\) 系数，就画有向关系 \(y\to x\)。令 \(x\preceq_L y\) 表示存在从 \(y\) 到 \(x\) 的有向链，允许空链。互相可达给出的等价类称为\term[zuo-baoqiang]{左胞腔}[left cell]。改用右乘定义右胞腔，同时允许两种乘法则定义双边胞腔。
\end{definition}
\symidx{\(x\preceq_L y\)：Hecke 规范基乘法定义的左胞腔预序}

这不是 Bruhat 偏序：它由乘法结构常数决定，互相可达还可能把不同长度的元素放在一起。定义采用的规范基版本必须一致；原文同时使用 \(C_w,C'_w\)，其等价的图定义及表示论用途见\cite[Section 1, Definition 1.2 and Theorems 1.3--1.4, pp. 166--167]{KazhdanLusztig1979Hecke}。

\ProseParagraphBreak
令 \(\Gamma\) 是一个左胞腔。所有满足 \(x\preceq_L y\)、某个 \(y\in\Gamma\) 的基元的线性生成空间是左理想，因为乘以 \(C'_s\) 只能沿有向关系继续向下走，而 \(H_s=C'_s-v^{-1}\)。其中不属于 \(\Gamma\) 的那些基元也生成左理想：若还能乘回 \(\Gamma\)，它就与 \(\Gamma\) 互相可达，矛盾。这两个左理想的商因而给出一个以 \(\Gamma\) 为基的胞腔表示。这一步只用定义与基的线性无关性。

\ProseParagraphBreak
在 \(S_3\) 中，置 \(b=v+v^{-1}\)，直接由前面的展开和二次关系得到
\[
(C'_s)^2=bC'_s,\quad
C'_tC'_s=C'_{ts},\quad
C'_sC'_{ts}=C'_{sts}+C'_s,\quad
C'_tC'_{ts}=bC'_{ts}.
\]
交换 \(s,t\) 给出另外四式，而 \(C'_sC'_{sts}=C'_tC'_{sts}=bC'_{sts}\)。由此左胞腔恰为
\[
\{1\},\qquad\{s,ts\},\qquad\{t,st\},\qquad\{sts\}.
\]
对 \(\Gamma=\{s,ts\}\)，上述两个左理想具体为
\[
I=R C'_s\oplus R C'_{ts}\oplus R C'_{sts},\qquad
J=R C'_{sts}.
\]
仅由 \(C'_s,C'_{ts}\) 张成的空间并不稳定，因为 \(C'_sC'_{ts}\) 含有 \(C'_{sts}\)。所取的是商模 \(I/J\)，令其中两基向量为 \(e_s=C'_s+J,e_{ts}=C'_{ts}+J\)。因此 \(C'_{sts}\) 项在商中为零，得到
\[
[C'_s]_\Gamma=
\begin{pmatrix}b&1\\0&0\end{pmatrix},
\qquad
[C'_t]_\Gamma=
\begin{pmatrix}0&0\\1&b\end{pmatrix}.
\]
由 \(H_i=C'_i-v^{-1}\) 得到 Hecke 生成元矩阵。在 \(v=1\) 处，它们成为
\[
s\longmapsto\begin{pmatrix}1&1\\0&-1\end{pmatrix},
\qquad
t\longmapsto\begin{pmatrix}-1&0\\1&1\end{pmatrix}.
\]
两者平方为一，乘积迹为 \(-1\)、行列式为一，且满足辫关系。三个换位的迹为零，两个三轮换的迹为 \(-1\)，所以该胞腔表示复化后与 \(S_3\) 的二维标准表示同构。胞腔方法由此把一个乘法基中的组合结构变成了实际的表示矩阵。
